%% kapa: exam=1 round=31 year=2020 period=1 form=A subject=civil_law subject_ko=민법 seq=40 number=40
%% source: Q-Net 감정평가사 기출문제 (cst00309) artlSeq=5207219 file=extracted/1008/5207219/A형/A형/2020년도 제31회 감정평가사 1차 시험_1교시_A형.pdf
%% conversion: kordoc 4.24.51 -> tools/gen1.py (auto); review=proofread-sonnet; flags=none
\begin{question}{40}
저당권에 관한 설명으로 옳지 않은 것은? (다툼이 있으면 판례에 따름)
\begin{choices}
  \item 저당권의 효력은 천연과실뿐만 아니라 법정과실에도 미친다.
  \item 저당권으로 담보된 채권을 양수하였으나 아직 대항요건을 갖추지 못한 양수인도 저당권이전의 부기등기를 마치고 저당권실행의 요건을 갖추면 경매를 신청할 수 있다.
  \item 후순위담보권자가 경매를 신청한 경우, 선순위근저당권의 피담보채권은 매수인이 경락대금을 완납하여 그 근저당권이 소멸하는 때에 확정된다.
  \item 저당권의 이전을 위하여 저당권의 양도인과 양수인, 그리고 저당권설정자 사이의 물권적 합의와 등기가 있어야 한다.
  \item 공동저당관계의 등기를 공동저당권의 성립요건이나 대항요건이라고는 할 수 없다.
\end{choices}
\end{question}
\answer{4}
